\chapter{Chemical Burns}

\section*{Definition and Overview}
A chemical burn is tissue damage caused by contact with a corrosive or reactive substance. Strong acids, alkalis, solvents, and oxidising agents can injure the skin, eyes, mouth, airway, or gastrointestinal tract depending on the route of exposure. The severity of a chemical burn depends on the nature and concentration of the chemical, the duration of contact, and the amount and area of exposure. Unlike thermal burns, chemical injury can continue to destroy tissue for minutes to hours after the initial contact unless the chemical is removed, so immediate and prolonged decontamination with water is the cornerstone of first aid. Chemical burns range from minor skin irritation to deep, full-thickness injuries and life-threatening systemic poisoning.

\section*{Epidemiology}
Chemical burns account for a small but important proportion of all burn injuries. They occur most often in occupational settings, particularly in manufacturing, agriculture, construction, and laboratories, and in industries that handle cleaning agents, acids, alkalis, and solvents. Household chemical burns are also common, usually from drain cleaners, oven cleaners, bleach, and disinfectants. Children are frequently injured by accidental exposure to household chemicals stored within reach. Chemical injuries to the eye are a leading cause of ocular trauma and can threaten vision if not treated immediately. Most chemical burns are preventable with appropriate handling, protective equipment, and safe storage.

\section*{Aetiology and Pathophysiology}
\begin{itemize}
    \item \textbf{Acids:} strong acids such as sulphuric, hydrochloric, and nitric acid denature proteins and cause coagulative necrosis, which tends to limit tissue penetration
    \item \textbf{Alkalis:} alkalis such as sodium and potassium hydroxide, ammonia, and cement cause liquefactive necrosis, dissolving tissue and penetrating deeply; they are often more dangerous than acids
    \item \textbf{Organic solvents:} substances such as petrol and industrial solvents dissolve fats and strip the skin's protective oils
    \item \textbf{Oxidising agents:} agents such as hydrogen peroxide and bleach damage tissue through oxidation and release of heat
    \item \textbf{Chemical warfare agents and phosphorus:} unusual exposures that cause specialised and often severe injuries
    \item \textbf{Mechanism of injury:} the chemical disrupts cellular proteins and membranes, causing cell death, tissue necrosis, and inflammation; damage continues while the agent remains in contact with tissue
    \item \textbf{Systemic toxicity:} some chemicals, such as hydrofluoric acid and phenol, are absorbed through the skin and can cause dangerous disturbances of calcium, heart rhythm, or other body systems
\end{itemize}

\section*{Clinical Features}
\begin{itemize}
    \item \textbf{Skin:} redness, swelling, pain, blistering, blackened or waxy-looking skin, and, in severe cases, full-thickness tissue loss
    \item \textbf{Eyes:} severe pain, redness, watering, photophobia, blurred vision, and in alkali injuries a characteristic "chalky white" cornea
    \item \textbf{Mouth and throat:} burns to the lips, tongue, and mouth from swallowing chemicals, with drooling and pain
    \item \textbf{Gastrointestinal tract:} abdominal pain, nausea, vomiting, and difficulty swallowing after ingestion
    \item \textbf{Airway:} hoarseness, stridor, coughing, and breathlessness if chemical fumes or hot liquids have been inhaled or aspirated
    \item \textbf{Systemic features:} confusion, drowsiness, low blood pressure, and abnormal heart rhythms from absorbed toxins
    \item Symptoms may be delayed, especially with alkali injuries, so the affected area should be inspected carefully and repeatedly
\end{itemize}

\section*{Differential Diagnosis}
\begin{itemize}
    \item Thermal burns, which have a similar appearance but a different mechanism
    \item Friction or contact dermatitis and allergic skin reactions
    \item Blistering skin disorders and infections such as staphylococcal scalded skin syndrome
    \item Chemical or drug-induced toxic epidermal necrolysis
    \item Conjunctivitis and corneal abrasion in eye injuries
    \item Oesophagitis, reflux, or other causes of sore throat and swallowing difficulty after ingestion
    \item Ingestion of other toxic or corrosive household substances
\end{itemize}

\section*{When to Seek Medical Care}
Any chemical burn to the eye, a burn that is deep, large, painful, or covers a significant area, a burn to the face or airway, or the ingestion of a corrosive chemical requires urgent medical assessment.

\medskip
\textbf{Call 000 (or local emergency services) for:}
\begin{itemize}
    \item Any chemical injury to the eye, especially from an alkali
    \item Difficulty breathing, hoarseness, or coughing after inhaling fumes or ingesting chemicals
    \item Swallowing of a corrosive chemical
    \item Fainting, confusion, seizures, or collapse after chemical exposure
    \item Extensive or deeply painful skin burns
\end{itemize}

\section*{Diagnosis and Investigations}
\begin{itemize}
    \item \textbf{First-aid history:} the identity of the chemical, its concentration, the route and duration of exposure, and whether the area was irrigated
    \item \textbf{Examination:} careful assessment of the skin, eyes, mouth, and airway, including checking the extent and depth of the burn
    \item \textbf{Eye examination:} fluorescein staining to identify corneal damage, with slit-lamp examination by an ophthalmologist for significant injuries
    \item \textbf{Laboratory tests:} blood gases, electrolytes, and tests for specific absorbed toxins (for example, calcium levels after hydrofluoric acid exposure)
    \item \textbf{Imaging:} chest X-ray or CT if fumes were inhaled or aspiration is suspected
    \item \textbf{Endoscopy:} examination of the oesophagus and stomach after ingestion of corrosives, performed carefully by a specialist to assess the extent of injury
    \item \textbf{Poisoning information:} the local poisons information service can provide specific guidance for the chemical involved
\end{itemize}

\section*{Management}
\begin{itemize}
    \item \textbf{Immediate decontamination:} remove contaminated clothing and rinse the affected area with copious running water for at least 20 minutes
    \item \textbf{Eye irrigation:} continuous irrigation of the eye with water or saline until the chemical is fully removed; the eye must take priority
    \item \textbf{Do not neutralise:} never add acid or alkali to "neutralise" a burn, as this generates heat and worsens the injury
    \item \textbf{Wound care:} cleaning, gentle debridement, and appropriate dressings, with referral to a burns service for deep or extensive injuries
    \item \textbf{Pain relief:} analgesics appropriate to the severity of the burn
    \item \textbf{Specialist treatment:} specific antidotes or washes for particular agents, such as calcium gluconate for hydrofluoric acid burns
    \item \textbf{Airway and breathing support:} oxygen, monitoring, and occasionally intubation for airway burns or severe inhalation injury
    \item \textbf{Gastrointestinal care:} specialist management after corrosive ingestion, which may include endoscopic assessment and rarely surgery for perforation
    \item \textbf{Tetanus prophylaxis:} update tetanus immunisation if due
    \item \textbf{Burns centre referral:} for deep, extensive, or cosmetically important burns, or those over joints, face, or perineum
\end{itemize}

\section*{Complications}
\begin{itemize}
    \item Deep skin necrosis and full-thickness burns requiring skin grafting
    \item Permanent scarring and contractures that limit movement
    \item Corneal ulceration, scarring, and permanent visual loss in eye injuries
    \item Oesophageal strictures and swallowing difficulty after corrosive ingestion
    \item Perforation of the stomach or oesophagus in severe ingestions
    \item Airway oedema and respiratory failure after inhalation injury
    \item Systemic poisoning, including low calcium, kidney injury, and cardiac arrest from absorbed chemicals
    \item Infection of the burn wound and sepsis
    \item Psychological distress, including anxiety and post-traumatic stress symptoms
\end{itemize}

\section*{Prognosis and Outcomes}
The outlook depends on the chemical involved, the depth and size of the burn, and how quickly effective decontamination was given. Superficial chemical burns typically heal within one to three weeks with good skin care and usually leave little scarring. Deep burns, especially alkali injuries, may leave permanent scarring, contractures, or functional loss. Eye injuries carry a risk of permanent visual impairment that increases with delayed irrigation and with alkali exposure. Outcomes are markedly better when copious washing begins immediately, so public awareness of correct first aid is vital. Most people with prompt treatment recover fully, and ongoing rehabilitation and physiotherapy support recovery after serious injuries.

\section*{Prevention and Self-Care}
\begin{itemize}
    \item Store household chemicals out of reach of children and in their original labelled containers
    \item Wear gloves, goggles, and protective clothing when handling acids, alkalis, bleaches, and cleaning agents
    \item Follow the manufacturer's safety instructions and use chemicals in well-ventilated areas
    \item Never mix cleaning products, as some combinations release toxic gases
    \item Keep a first-aid supply of water close by when working with chemicals
    \item Rinse any accidental skin contact with running water immediately
    \item If a chemical is splashed into the eye, irrigate with water immediately and seek urgent medical care
    \item In the workplace, use the mandated personal protective equipment and follow hazardous-substances procedures
\end{itemize}

\section*{Sources and Further Reading}
The following reputable sources provide further information on chemical burns:
\begin{itemize}
    \item NHS. \textit{Health A--Z: burns and scalds information.} https://www.nhs.uk/conditions/
    \item Mayo Clinic. \textit{First aid and emergency information.} https://www.mayoclinic.org/
    \item MSD Manual. \textit{Consumer and professional editions on burns and poisoning.} https://www.msdmanuals.com/
    \item MedlinePlus. \textit{Health topics on burns and safety.} https://medlineplus.gov/
    \item World Health Organization. \textit{Burn prevention and management fact sheets.} https://www.who.int/
\end{itemize}
